% ============================================================
%  CHAPTER 28: RADIOACTIVITY AND NUCLEAR PHYSICS
%  The Physics Codex: A Complete University Foundation
%  Korede Samuel Omotosho (Falcon)
%  Aurikrex Learning Systems, 2026
%
%  File: chapters/part5/ch28_radioactivity.tex
% ============================================================

\chapter{Radioactivity and Nuclear Physics}
\label{ch:radioactivity}
\minitoc
\newpage

\chapterquote{We are all made of stardust. The heavy
elements inside you --- the iron in your blood, the
calcium in your bones --- were forged in the cores of
stars that exploded before our solar system was born.}{Carl Sagan}

\begin{objectivebox}
By the end of this chapter, you should be able to:
\begin{enumerate}[noitemsep]
  \item Describe the structure of the nucleus using
  proton and neutron numbers
  \item Describe the three types of radioactive
  emission and their properties
  \item Write and balance nuclear equations for
  alpha decay, beta decay and gamma emission
  \item Explain radioactive decay as a random,
  spontaneous process
  \item Define and apply the concept of half-life
  and the decay constant
  \item Apply the exponential decay law
  \item Define and apply nuclear binding energy
  and the mass-energy equivalence $E = mc^2$
  \item Describe and distinguish between nuclear
  fission and fusion
  \item Describe the operation of a nuclear reactor
  \item Apply knowledge of radioactivity to
  practical uses including medical imaging, carbon
  dating and radiation safety
\end{enumerate}
\end{objectivebox}

% ============================================================
\section{Intuition: Energy from the Nucleus}
% ============================================================

The Sun converts $600$ million tonnes of hydrogen into
helium \textit{every second}. In doing so, it releases
energy equivalent to the annihilation of 4 million
tonnes of matter per second via Einstein's formula
$E = mc^2$. This energy has sustained life on Earth
for 4.6 billion years and will continue for another
5 billion.

On Earth, we have discovered how to release nuclear
energy in controlled and uncontrolled ways. Nuclear
weapons release energy in microseconds. Nuclear power
stations release it over months and years. Both exploit
the same fundamental physics: the extraordinarily large
binding energy per nucleon in atomic nuclei.

Radioactivity --- the spontaneous emission of radiation
from unstable nuclei --- is the natural version of the
same process. Understanding it is essential for medicine,
energy generation, archaeology, and safety.

% ============================================================
\section{Nuclear Structure}
% ============================================================

\begin{definitionbox}[Nuclear Terminology]
The \textbf{nucleus} of an atom contains
\textbf{protons} and \textbf{neutrons}, collectively
called \textbf{nucleons}.

\textbf{Proton number} (atomic number) $Z$: the
number of protons in the nucleus. Determines the
element (e.g.\ $Z = 6$ is always carbon).

\textbf{Nucleon number} (mass number) $A$: the total
number of protons plus neutrons.

\textbf{Neutron number} $N = A - Z$.

Standard notation:
\[
  {}^{A}_{Z}\text{X}
\]
where X is the chemical symbol.

\textbf{Isotopes:} atoms of the same element (same $Z$)
with different numbers of neutrons (different $A$).
Example: ${}^{12}_{6}\text{C}$, ${}^{13}_{6}\text{C}$,
${}^{14}_{6}\text{C}$ are all isotopes of carbon.

\textbf{Nuclear radius:} approximately
$R = R_0 A^{1/3}$ where $R_0 \approx 1.2\times10^{-15}\ \text{m}$.
Nuclear density is approximately constant for all
nuclei: $\sim 10^{17}\ \text{kg\,m}^{-3}$.
\end{definitionbox}

\begin{keybox}[The Four Fundamental Forces in the Nucleus]
The nucleus is incredibly dense. Protons are packed
together despite their mutual electrostatic repulsion.
What holds them together?

The \textbf{strong nuclear force} --- one of the four
fundamental forces. It is:
\begin{itemize}[noitemsep]
  \item About 100 times stronger than the
  electromagnetic force
  \item Very short range ($\sim 10^{-15}\ \text{m}$
  --- only acts between neighbouring nucleons)
  \item Acts equally between proton-proton,
  neutron-neutron, and proton-neutron pairs
  \item Attractive at short range; repulsive at
  even shorter range (prevents nucleons merging)
\end{itemize}
For large nuclei, the long-range electrostatic
repulsion between protons begins to overwhelm the
short-range strong force, making heavy nuclei
increasingly unstable.
\end{keybox}

% ============================================================
\section{Radioactive Emissions}
% ============================================================

Some nuclei are unstable and spontaneously emit
radiation to reach a more stable configuration. This
process is \textbf{radioactive decay}. There are
three main types of emission.

\begin{table}[H]
\centering
\caption{Properties of Alpha, Beta and Gamma Radiation}
\label{tab:radiation_types}
\small
\begin{tabular}{llllll}
\toprule
\textbf{Type} & \textbf{Symbol} &
\textbf{Composition} & \textbf{Charge} &
\textbf{Mass} & \textbf{Penetration}\\
\midrule
Alpha & $\alpha$ & 2 protons + 2 neutrons (${}^4_2\text{He}$) &
  $+2e$ & $\sim 4$ u & Few cm air; stopped by paper\\
Beta-minus & $\beta^-$ & Fast electron & $-e$ &
  $\sim 0$ & Few mm aluminium\\
Beta-plus & $\beta^+$ & Fast positron & $+e$ &
  $\sim 0$ & Few mm aluminium\\
Gamma & $\gamma$ & High-energy photon & 0 &
  0 & Several cm lead; never fully stopped\\
\bottomrule
\end{tabular}
\end{table}

\begin{table}[H]
\centering
\caption{Ionising Power and Deflection in Fields}
\label{tab:radiation_properties}
\small
\begin{tabular}{llll}
\toprule
\textbf{Type} & \textbf{Ionising power} &
\textbf{In electric field} & \textbf{In magnetic field}\\
\midrule
Alpha & Strongly ionising & Deflected toward $-$ plate &
  Deflected; large radius\\
Beta-minus & Moderately ionising & Deflected toward $+$ plate &
  Deflected; small radius, opposite to $\alpha$\\
Gamma & Weakly ionising & Not deflected &
  Not deflected\\
\bottomrule
\end{tabular}
\end{table}

\begin{diagbox}[Alpha, Beta and Gamma in Electric and Magnetic Fields]
\begin{center}
\begin{tikzpicture}[scale=0.9, font=\small]

  % === ELECTRIC FIELD DEFLECTION ===
  \begin{scope}[shift={(0,0)}]
    \node[font=\bfseries] at (4.0,6.5)
      {Deflection in Electric Field};

    % Plates
    \fill[codexred!30] (1.5,5.5) rectangle (6.5,6.0);
    \fill[codexblue!30] (1.5,0.5) rectangle (6.5,1.0);
    \draw[thick] (1.5,5.5) -- (6.5,5.5);
    \draw[thick] (1.5,1.0) -- (6.5,1.0);
    \node[right, codexred] at (6.6,5.75) {$+$};
    \node[right, codexblue] at (6.6,0.75) {$-$};

    % Source
    \fill[gray!50] (0.5,2.8) rectangle (1.2,3.7);
    \draw[thick] (0.5,2.8) rectangle (1.2,3.7);
    \node[left, font=\tiny] at (0.4,3.25)
      {Source};

    % Radiation beams
    % Positive alpha particles deflect toward the negative lower plate.
    \draw[->, very thick, codexred]
      (1.2,3.25) -- (6.0,2.0);
    \node[below, codexred] at (6.1,2.0)
      {$\alpha$ (toward $-$ plate)};

    % Gamma (straight through)
    \draw[->, very thick, codexgold]
      (1.2,3.25) -- (6.5,3.25);
    \node[right, codexgold] at (6.6,3.25)
      {$\gamma$ (undeflected)};

    % Negative beta particles deflect toward the positive upper plate.
    \draw[->, very thick, codexblue]
      (1.2,3.25) -- (6.0,4.5);
    \node[above, codexblue] at (6.1,4.5)
      {$\beta^-$ (toward $+$ plate)};

    % Electric field arrows
    \foreach \x in {2,3,4,5,6} {
      \draw[->, gray!50] (\x,5.5) -- (\x,1.0);
    }
    \node[codexred, font=\tiny] at (2.0,5.3)
      {$\vect{E}$};
  \end{scope}

  % === MAGNETIC FIELD DEFLECTION ===
  \begin{scope}[shift={(9,0)}]
    \node[font=\bfseries, align=center] at (4.0,6.5)
      {Deflection in Magnetic Field\\($\vect{B}$ into page)};

    % Field region (crosses)
    \fill[gray!10] (1.5,1.0) rectangle (6.5,5.5);
    \draw[thick] (1.5,1.0) rectangle (6.5,5.5);
    \foreach \x in {2,3,4,5,6} {
      \foreach \y in {1.5,2.5,3.5,4.5} {
        \node[gray!50, font=\small] at (\x,\y) {$\times$};
      }
    }
    \node[gray, font=\tiny] at (2.0,5.2)
      {$\vect{B}$ into page};

    % Source
    \fill[gray!50] (0.5,2.8) rectangle (1.2,3.7);
    \draw[thick] (0.5,2.8) rectangle (1.2,3.7);

    % Alpha bends gently upward (large radius)
    \draw[->, very thick, codexred]
      (1.2,3.25) .. controls (3.0,3.25) and (5.0,3.55)
      .. (6.3,4.15);
    \node[above, codexred, font=\tiny] at (4.0,3.65)
      {$\alpha$ (large radius)};

    % Gamma (straight)
    \draw[->, very thick, codexgold]
      (1.2,3.25) -- (6.5,3.25);
    \node[above, codexgold] at (5.0,3.35) {$\gamma$};

    % Beta-minus bends sharply downward (small radius, opposite charge)
    \draw[->, very thick, codexblue]
      (1.2,3.25) .. controls (1.6,3.25) and (2.1,2.55)
      .. (2.6,1.75);
    \node[below, codexblue, font=\tiny] at (2.2,1.55)
      {$\beta^-$ (small radius)};

    \node[align=center, font=\tiny, gray]
      at (4.0,0.5)
      {$\alpha$ and $\beta$ curve in opposite directions\\
      (opposite charges). $\gamma$ is unaffected.};
  \end{scope}

\end{tikzpicture}
\end{center}
\end{diagbox}

% ============================================================
\section{Nuclear Decay Equations}
% ============================================================

\begin{lawbox}[Conservation Laws in Nuclear Decay]
In every nuclear decay:
\begin{itemize}[noitemsep]
  \item Total \textbf{nucleon number} $A$ is
  conserved (sum of mass numbers is constant).
  \item The nuclear proton number $Z$ can change by one
  in beta decay.
  \item Total electric charge, including emitted beta
  particles, is conserved.
  \item Total \textbf{energy-momentum} is conserved.
\end{itemize}
\end{lawbox}

\begin{definitionbox}[Alpha Decay]
Alpha decay: the nucleus emits an alpha particle
(${}^4_2\text{He}$). The parent nucleus loses 2
protons and 2 neutrons:
\[
  {}^A_Z\text{X} \rightarrow {}^{A-4}_{Z-2}\text{Y}
  + {}^4_2\text{He}
\]
Example: Radium-226 decays to Radon-222:
\[
  {}^{226}_{88}\text{Ra} \rightarrow
  {}^{222}_{86}\text{Rn} + {}^{4}_{2}\text{He}
\]
Check: $A$: $226 = 222 + 4$ ✓ \quad $Z$: $88 = 86 + 2$ ✓

Alpha emission typically occurs in heavy nuclei
($A > 200$) where the nucleus is too large to be
held together by the short-range strong force.
\end{definitionbox}

\begin{definitionbox}[Beta-Minus Decay]
Beta-minus ($\beta^-$) decay: a neutron in the nucleus
converts to a proton, emitting a fast electron and an
antineutrino $\bar{\nu}_e$:
\[
  n \rightarrow p + e^- + \bar{\nu}_e
\]
The nucleus gains one proton and loses one neutron:
\[
  {}^A_Z\text{X} \rightarrow {}^{A}_{Z+1}\text{Y}
  + {}^{0}_{-1}e + \bar{\nu}_e
\]
Example: Carbon-14 decays to Nitrogen-14:
\[
  {}^{14}_{6}\text{C} \rightarrow
  {}^{14}_{7}\text{N} + {}^{0}_{-1}e + \bar{\nu}_e
\]
Beta-minus occurs when the nucleus has too many
neutrons relative to protons.
\end{definitionbox}

\begin{definitionbox}[Beta-Plus Decay and Gamma Emission]
\textbf{Beta-plus} ($\beta^+$) decay: a proton converts
to a neutron, emitting a positron and neutrino:
\[
  {}^A_Z\text{X} \rightarrow {}^{A}_{Z-1}\text{Y}
  + {}^{0}_{+1}e + \nu_e
\]
Occurs when nucleus has too many protons.

\textbf{Gamma emission:} A daughter nucleus produced by
alpha or beta decay may be left in an excited state. It
can release excess energy as a gamma photon, without
changing $A$ or $Z$:
\[
  {}^A_Z\text{X}^* \rightarrow {}^A_Z\text{X} + \gamma
\]
\end{definitionbox}

\begin{examplebox}[28.1]
\stepcounter{workedexample}
Complete the following nuclear equations:

(a) ${}^{238}_{92}\text{U} \rightarrow {}^{?}_{?}\text{Th}
+ {}^{4}_{2}\text{He}$

(b) ${}^{131}_{53}\text{I} \rightarrow {}^{?}_{?}\text{Xe}
+ {}^{0}_{-1}e + \bar{\nu}_e$

(c) ${}^{22}_{11}\text{Na} \rightarrow {}^{?}_{?}\text{Ne}
+ {}^{0}_{+1}e + \nu_e$
\end{examplebox}

\begin{solutionbox}
\textbf{(a)} Alpha decay: $A: 238-4=234$;
$Z: 92-2=90$ (Thorium):
\[
  {}^{238}_{92}\text{U} \rightarrow
  {}^{234}_{90}\text{Th} + {}^{4}_{2}\text{He}
\]

\textbf{(b)} Beta-minus: $A: 131$ (unchanged);
$Z: 53+1=54$ (Xenon):
\[
  {}^{131}_{53}\text{I} \rightarrow
  {}^{131}_{54}\text{Xe} + {}^{0}_{-1}e
  + \bar{\nu}_e
\]

\textbf{(c)} Beta-plus: $A: 22$ (unchanged);
$Z: 11-1=10$ (Neon):
\[
  {}^{22}_{11}\text{Na} \rightarrow
  {}^{22}_{10}\text{Ne} + {}^{0}_{+1}e + \nu_e
\]
\end{solutionbox}

% ============================================================
\section{Radioactive Decay: The Statistical Nature}
% ============================================================

\begin{defbox}[Radioactive Decay: Random and Spontaneous]
Radioactive decay is:
\begin{itemize}[noitemsep]
  \item \textbf{Spontaneous:} It cannot be triggered
  or prevented by chemical, physical or environmental
  conditions (temperature, pressure, chemical state).
  \item \textbf{Random:} It is impossible to predict
  which particular nucleus in a sample will decay next
  or when. Only statistical averages are meaningful.
\end{itemize}

The \textbf{decay constant} $\lambda$ is the
probability of a given nucleus decaying per unit time.
For a sample of $N$ nuclei, the \textbf{activity}
$A$ (decays per second) is:
\[
  A = -\frac{dN}{dt} = \lambda N
\]
SI unit of activity: becquerel (Bq),
where $1\ \text{Bq} = 1\ \text{decay per second}$.
Older unit: curie (Ci), $1\ \text{Ci}
= 3.7\times10^{10}\ \text{Bq}$.
\end{defbox}

% ============================================================
\section{Half-Life and Exponential Decay}
% ============================================================

\begin{definitionbox}[Half-Life]
The \textbf{half-life} $t_{1/2}$ of a radioactive
isotope is the time for:
\begin{itemize}[noitemsep]
  \item Half of the original number of nuclei to
  decay, OR
  \item The activity of the sample to halve, OR
  \item The number of undecayed nuclei to halve.
\end{itemize}

Relationship between half-life and decay constant:
\[
  t_{1/2} = \frac{\ln 2}{\lambda} = \frac{0.693}{\lambda}
\]
\end{definitionbox}

\begin{lawbox}[Exponential Decay Law]
The number of undecayed nuclei at time $t$:
\[
  N = N_0\,e^{-\lambda t}
  = N_0\left(\frac{1}{2}\right)^{t/t_{1/2}}
\]
The activity at time $t$:
\[
  A = A_0\,e^{-\lambda t}
  = \lambda N_0\,e^{-\lambda t}
\]
where $N_0$ is the initial number of nuclei and
$A_0 = \lambda N_0$ is the initial activity.
\end{lawbox}

\begin{diagbox}[Radioactive Decay: Exponential Decay Curve and Half-Lives]
\begin{center}
\begin{tikzpicture}
\begin{axis}[
  width=13cm, height=7cm,
  xlabel={Time $t$ (in half-lives)},
  ylabel={Fraction remaining ($N/N_0$)},
  xmin=0, xmax=5,
  ymin=0, ymax=1.2,
  grid=major, grid style={gray!15},
  xtick={0,1,2,3,4,5},
  xticklabels={0,$t_{1/2}$,$2t_{1/2}$,$3t_{1/2}$,
    $4t_{1/2}$,$5t_{1/2}$},
  ytick={0,0.0625,0.125,0.25,0.5,1.0},
  yticklabels={0,$1/16$,$1/8$,$1/4$,$1/2$,1},
]
  % Exponential decay curve
  \addplot[codexblue, very thick, domain=0:5,
    samples=100]
    {exp(-0.693*x)};

  % Horizontal dashed lines at each half-life fraction
  \addplot[dashed, codexred!50, domain=0:1]{0.5};
  \addplot[dashed, codexred!50, domain=0:2]{0.25};
  \addplot[dashed, codexred!50, domain=0:3]{0.125};
  \addplot[dashed, codexred!50, domain=0:4]{0.0625};

  % Vertical dashed lines
  \draw[dashed, codexred!50]
    (axis cs:1,0) -- (axis cs:1,0.5);
  \draw[dashed, codexred!50]
    (axis cs:2,0) -- (axis cs:2,0.25);
  \draw[dashed, codexred!50]
    (axis cs:3,0) -- (axis cs:3,0.125);
  \draw[dashed, codexred!50]
    (axis cs:4,0) -- (axis cs:4,0.0625);

  % Data points (use a pgfplots-native plot to avoid macro-scope issues)
  \addplot[only marks, mark=*, mark size=2.2pt, codexblue]
    coordinates {(0,1) (1,0.5) (2,0.25) (3,0.125) (4,0.0625) (5,0.03125)};

  % Annotations
  \node[above right, codexblue, font=\small]
    at (axis cs:0,1.0) {$N_0$};
  \node[right, codexblue, font=\small]
    at (axis cs:1,0.52) {$N_0/2$};
  \node[right, codexblue, font=\small]
    at (axis cs:2,0.27) {$N_0/4$};
  \node[right, codexblue, font=\small]
    at (axis cs:3,0.14) {$N_0/8$};

  % Table in corner
  \node[align=left, font=\tiny, fill=white,
    draw=gray, inner sep=4pt]
    at (axis cs:4,0.7)
    {After $n$ half-lives:\\$N = N_0(1/2)^n$};

\end{axis}
\end{tikzpicture}
\end{center}
\end{diagbox}

\begin{examplebox}[28.2]
\stepcounter{workedexample}
Iodine-131 has a half-life of $8.0$ days. A hospital
receives a sample with initial activity
$4.0\times10^8\ \text{Bq}$.\\
(a) Calculate the decay constant.\\
(b) Calculate the activity after $24$ days.\\
(c) Calculate the number of ${}^{131}$I atoms
initially present.\\
(d) How long before the activity falls to
$5.0\times10^6\ \text{Bq}$?
\end{examplebox}

\begin{solutionbox}
$t_{1/2} = 8.0\ \text{days} = 8.0 \times 86400
= 6.912\times10^5\ \text{s}$

\textbf{(a)} Decay constant:
\[
  \lambda = \frac{\ln 2}{t_{1/2}}
  = \frac{0.693}{6.912\times10^5}
  = 1.003\times10^{-6}\ \text{s}^{-1}
\]

\textbf{(b)} After $24$ days $= 3 \times t_{1/2}$:
\[
  A = A_0\left(\frac{1}{2}\right)^3
  = 4.0\times10^8 \times \frac{1}{8}
  = 5.0\times10^7\ \text{Bq}
\]

\textbf{(c)} Initial number of atoms:
\[
  N_0 = \frac{A_0}{\lambda}
  = \frac{4.0\times10^8}{1.003\times10^{-6}}
  = 3.99\times10^{14}\ \text{atoms}
  \approx 4.0\times10^{14}\ \text{atoms}
\]

\textbf{(d)} Time to reach $5.0\times10^6\ \text{Bq}$:
\[
  \frac{A}{A_0} = \frac{5.0\times10^6}{4.0\times10^8}
  = 0.0125 = \left(\frac{1}{2}\right)^n
\]
\[
  n = \frac{\ln(0.0125)}{\ln(0.5)}
  = \frac{-4.382}{-0.693} = 6.32\ \text{half-lives}
\]
\[
  t = 6.32 \times 8.0 = 50.6\ \text{days}
\]
\end{solutionbox}

% ============================================================
\section{Applications of Radioactivity}
% ============================================================

\subsection{Radiocarbon Dating}

\begin{processbox}[Carbon-14 Dating]
${}^{14}$C is continuously produced in the upper
atmosphere by cosmic ray neutrons colliding with
nitrogen:
\[
  {}^{14}_{7}\text{N} + {}^{1}_{0}\text{n}
  \rightarrow {}^{14}_{6}\text{C}
  + {}^{1}_{1}\text{H}
\]
${}^{14}$C (half-life $= 5730$ years) is absorbed
by living organisms from atmospheric CO$_2$. While
alive, organisms maintain a constant ratio of
${}^{14}$C to ${}^{12}$C in equilibrium with the
atmosphere. When they die, no more ${}^{14}$C is
absorbed and the existing ${}^{14}$C decays.

By measuring the current activity of a sample and
comparing with the expected activity for a living
organism, the age can be calculated:
\[
  t = \frac{t_{1/2}}{\ln 2}\ln\left(\frac{A_0}{A}\right)
\]
Reliable for organic material up to about $50{,}000$
years old.
\end{processbox}

\subsection{Medical Applications}

\begin{table}[H]
\centering
\caption{Medical Uses of Radioactive Isotopes}
\label{tab:medical_radioactivity}
\begin{tabular}{llll}
\toprule
\textbf{Isotope} & \textbf{Type} & \textbf{Half-life} &
\textbf{Medical Use}\\
\midrule
Technetium-99m & $\gamma$ & 6 hours & Diagnostic imaging
(SPECT scans)\\
Iodine-131 & $\beta^-$, $\gamma$ & 8 days &
  Thyroid cancer treatment\\
Fluorine-18 & $\beta^+$ & 110 min & PET scans\\
Cobalt-60 & $\gamma$ & 5.27 years & Radiotherapy\\
Iridium-192 & $\beta^-$, $\gamma$ & 74 days &
  Brachytherapy (cancer)\\
\bottomrule
\end{tabular}
\end{table}

% ============================================================
\section{Nuclear Binding Energy and \texorpdfstring{$E = mc^2$}{E = mc2}}
% ============================================================

\begin{definitionbox}[Mass-Energy Equivalence]
Einstein's special theory of relativity (1905) shows
that mass and energy are interconvertible:
\[
  E = mc^2
\]
where $c = 3.0\times10^8\ \text{m\,s}^{-1}$.

The \textbf{atomic mass unit} (u):
$1\ \text{u} = 1.661\times10^{-27}\ \text{kg}$,
which corresponds to energy:
\[
  1\ \text{u} \times c^2 = 931.5\ \text{MeV}
\]

Useful masses for nuclear calculations:
\begin{itemize}[noitemsep]
  \item Proton: $m_p = 1.007276\ \text{u}$
  \item Neutron: $m_n = 1.008665\ \text{u}$
  \item Electron: $m_e = 0.000549\ \text{u}$
  \item Neutral hydrogen-1 atom: $m({}^1_1\text{H}) = 1.007825\ \text{u}$
  \item Neutral helium-4 atom: $m({}^4_2\text{He}) = 4.002603\ \text{u}$
\end{itemize}
\end{definitionbox}

\begin{definitionbox}[Nuclear Binding Energy]
Using nuclear masses consistently, the \textbf{mass defect}
$\Delta m$ is the sum of the constituent nucleon masses
minus the nuclear mass:
\[
  \Delta m = Zm_p + (A-Z)m_n - m_{\mathrm{nucleus}}
\]
If neutral atomic masses are used instead, use
\[
  \Delta m = Zm({}^1_1\text{H}) + (A-Z)m_n - m_{\mathrm{atom}}
\]
where $m({}^1_1\text{H})$ is the neutral hydrogen-1 atom
mass and $m_{\mathrm{atom}}$ is the neutral nuclide atom
mass. The electron masses then cancel, apart from small
differences in electron-binding energies.
The \textbf{binding energy} is the energy equivalent
of this mass defect:
\[
  E_B = \Delta m \cdot c^2
\]
The \textbf{binding energy per nucleon} $E_B/A$ is a
measure of nuclear stability. It peaks near iron
(${}^{56}\text{Fe}$) at about $8.8\ \text{MeV}$ per
nucleon, which is therefore the most stable nucleus.

\textbf{Physical meaning:} The binding energy is the
energy that must be supplied to completely separate
all the nucleons in a nucleus. Equivalently, it is
the energy released when those nucleons come together
to form the nucleus from scratch.
\end{definitionbox}

\begin{diagbox}[Binding Energy Per Nucleon vs Mass Number]
\begin{center}
\begin{tikzpicture}
\begin{axis}[
  width=13cm, height=6.5cm,
  xlabel={Mass number $A$ (nucleon number)},
  ylabel={Binding energy per nucleon (MeV)},
  xmin=0, xmax=240,
  ymin=0, ymax=10,
  grid=major, grid style={gray!15},
  xtick={0,20,40,60,80,100,120,140,160,180,200,220,240},
  ytick={0,1,2,3,4,5,6,7,8,9,10},
]
  % Approximate binding energy curve
  % Rises steeply from H, peaks near Fe (A=56)
  % then slowly falls for heavy elements

  % Low-A points (including the zero binding of hydrogen-1)
  \addplot[codexblue, very thick]
    coordinates {(1,0) (2,1.11) (3,2.57) (4,7.07)};
  % H-1 at 0
  \fill[codexblue] (axis cs:1,0) circle (3pt);
  \node[above, font=\tiny] at (axis cs:1,0.2)
    {$^1$H (0)};
  % He-4 at 7.07
  \fill[codexblue] (axis cs:4,7.07) circle (3pt);
  \node[above, font=\tiny] at (axis cs:6,7.2)
    {$^4$He};

  % Rise through helium-4 and carbon-12 to the iron-56 peak
  \addplot[codexblue, very thick, domain=4:56,
    samples=50]
    {8.8 - (8.8-7.07)*((56-x)/52)^2.6027};

  % Peak at Fe-56
  \fill[codexred] (axis cs:56,8.8) circle (5pt);
  \node[above, codexred, font=\small]
    at (axis cs:56,8.85) {${}^{56}$Fe (most stable)};

  % Slow decline for heavy nuclei
  \addplot[codexblue, very thick, domain=56:240,
    samples=100]
    {8.8 - 1.23*(x-56)/182};

  % Key nuclei marked
  \fill[codexgold] (axis cs:12,7.68) circle (3pt);
  \node[below, codexgold, font=\tiny]
    at (axis cs:12,7.6) {${}^{12}$C};
  \fill[codexgold] (axis cs:238,7.57) circle (3pt);
  \node[below, codexgold, font=\tiny]
    at (axis cs:228,7.47) {${}^{238}$U};

  % Fission annotation
  \draw[->, thick, codexred, dashed]
    (axis cs:235,7.61) -- (axis cs:120,8.4);
  \node[above, codexred, font=\small]
    at (axis cs:160,8.1)
    {Fission: heavy $\to$ medium\\releases energy};

  % Fusion annotation
  \draw[->, thick, codexblue, dashed]
    (axis cs:2,1.1) -- (axis cs:20,7.4);
  \node[below, codexblue, font=\small, align=center]
    at (axis cs:70,5.5)
    {Fusion: light $\to$ medium\\releases energy};

\end{axis}
\end{tikzpicture}
\end{center}
\small Energy is released both by fusing light nuclei
(moving up the curve from the left) and by splitting
heavy nuclei (moving up from the right). In both cases,
the products have higher binding energy per nucleon
than the reactants --- the ``extra'' binding energy
is released as kinetic energy and radiation.
\end{diagbox}

\begin{examplebox}[28.3]
\stepcounter{workedexample}
Calculate the binding energy of helium-4
(${}^4_2\text{He}$) in MeV.\\
Use neutral-atom masses:
$m({}^1_1\text{H}) = 1.007825\ \text{u}$,
$m_n = 1.008665\ \text{u}$, and
$m({}^4_2\text{He}) = 4.002603\ \text{u}$.
$1\ \text{u} = 931.5\ \text{MeV}/c^2$.
\end{examplebox}

\begin{solutionbox}
Helium-4 has 2 protons and 2 neutrons.

Mass of constituents, using the neutral-atomic-mass form:
\[
  m_{parts} = 2m({}^1_1\text{H}) + 2m_n
  = 2(1.007825) + 2(1.008665)
  = 2.015650 + 2.017330
  = 4.032980\ \text{u}
\]

Mass defect:
\[
  \Delta m = m_{parts} - m({}^4_2\text{He})
  = 4.032980 - 4.002603
  = 0.030377\ \text{u}
\]

Binding energy:
\[
  E_B = \Delta m \times 931.5\ \text{MeV/u}
  = 0.030377 \times 931.5
  = 28.296\ \text{MeV}
\]

Binding energy per nucleon:
\[
  E_B/A = 28.296/4 = 7.074\ \text{MeV per nucleon}
\]

This is significantly lower than iron ($8.8\ \text{MeV}$),
consistent with the curve. Helium-4 is stable but
not the most stable nucleus.
\end{solutionbox}

% ============================================================
\section{Nuclear Fission}
% ============================================================

\begin{definitionbox}[Nuclear Fission]
\textbf{Nuclear fission} is the splitting of a heavy
nucleus into two lighter nuclei (fission fragments)
of roughly equal mass, releasing energy and typically
2--3 neutrons.

Example: thermal neutron-induced fission of Uranium-235:
\[
  {}^{1}_{0}\text{n} + {}^{235}_{92}\text{U}
  \rightarrow {}^{236}_{92}\text{U}^*
  \rightarrow {}^{141}_{56}\text{Ba}
  + {}^{92}_{36}\text{Kr}
  + 3{}^{1}_{0}\text{n}
\]
Energy released per fission: $\sim 200\ \text{MeV}$
(enormous compared to $\sim 1\ \text{eV}$ per chemical
reaction).

\textbf{Chain reaction:} The 2--3 neutrons released
can each trigger further fission events. If on average
more than one neutron from each fission causes another
fission (the multiplication factor $k > 1$), the
reaction is supercritical and the number of fissions
grows exponentially --- this is a nuclear explosion.

\textbf{Critical mass:} The minimum mass of fissile
material needed to sustain a self-sustaining chain
reaction ($k = 1$).
\end{definitionbox}

% ============================================================
\section{The Nuclear Reactor}
% ============================================================

\begin{processbox}[How a Nuclear Fission Reactor Works]
In a controlled fission reactor ($k = 1$):

\textbf{Fuel:} Enriched uranium (${}^{235}$U,
typically 3--5\% for power reactors).

\textbf{Moderator:} Slows down fast neutrons to
thermal speeds (slow neutrons are more effective
at causing fission). Common moderators: water,
heavy water (D$_2$O), graphite.

\textbf{Control rods:} Absorb neutrons to control
the reaction rate. Made of boron or cadmium. Inserted
deeper $\Rightarrow$ fewer neutrons $\Rightarrow$
slower reaction. This is how reactor power is
regulated.

\textbf{Coolant:} Carries heat from the reactor core
to a heat exchanger (steam generator). Common
coolants: water, CO$_2$, helium.

\textbf{Heat exchanger:} Transfers heat to produce
steam to drive turbines and generators.

\textbf{Containment:} Thick concrete and steel
shielding to absorb neutrons and gamma radiation.
\end{processbox}

\begin{diagbox}[Schematic of a Pressurised Water Reactor (PWR)]
\begin{center}
\begin{tikzpicture}[scale=0.85, font=\small]

  % === Reactor vessel ===
  \fill[gray!20] (0,0) rectangle (4.5,6.5);
  \draw[very thick] (0,0) rectangle (4.5,6.5);
  \node[font=\bfseries] at (2.25,7.0)
    {Reactor Vessel};

  % Fuel rods
  \foreach \x in {0.6,1.2,1.8,2.4,3.0,3.6} {
    \fill[codexgold!60] (\x,0.5) rectangle (\x+0.3,5.5);
    \draw[thick, codexgold] (\x,0.5) rectangle
      (\x+0.3,5.5);
  }
  % Control rods (partially inserted from top)
  \foreach \x in {0.9,1.5,2.1,2.7,3.3} {
    \fill[codexblue!60] (\x,3.5) rectangle
      (\x+0.2,6.5);
    \draw[thick, codexblue] (\x,3.5) rectangle
      (\x+0.2,6.5);
  }
  % Coolant (water as moderator+coolant)
  \fill[codexblue!15] (0.3,0.3) rectangle (4.2,6.2);
  % Fuel rods on top
  \foreach \x in {0.6,1.2,1.8,2.4,3.0,3.6} {
    \fill[codexgold!60] (\x,0.5) rectangle
      (\x+0.3,5.5);
    \draw[thick, codexgold] (\x,0.5) rectangle
      (\x+0.3,5.5);
  }
  % Compact key below the vessel keeps component labels clear of
  % the primary-coolant arrows in the reactor--exchanger gap.
  \node[draw=gray!50, fill=white, rounded corners, inner sep=2pt,
    text width=3.4cm, align=center, font=\tiny] at (2.25,-0.75)
    {Fuel rods (${}^{235}$U); control rods; pressurised water:
     moderator + coolant};

  % Hot coolant out
  \draw[->, very thick, codexred]
    (4.5,4.5) -- (6.5,4.5);
  \node[above, codexred, font=\tiny] at (5.5,4.9)
    {Hot coolant};

  % Cool coolant in
  \draw[->, very thick, codexblue]
    (6.5,1.5) -- (4.5,1.5);
  \node[below, codexblue, font=\tiny] at (5.5,1.1)
    {Cool coolant};

  % === Heat exchanger ===
  \fill[gray!15] (6.5,0.5) rectangle (9.5,6.0);
  \draw[thick] (6.5,0.5) rectangle (9.5,6.0);
  \node[font=\bfseries] at (8.0,6.5)
    {Heat Exchanger};
  % Separate primary-coolant and secondary water/steam paths.
  % Primary hot coolant enters from the reactor and returns cooled.
  \draw[->, very thick, codexred]
    (6.5,4.5) -- (7.2,4.5)
    .. controls (7.5,4.4) and (7.5,3.8) .. (7.5,3.0);
  \draw[->, very thick, codexblue]
    (7.5,3.0) .. controls (7.5,2.2) and (7.5,1.8) .. (7.2,1.5)
    -- (6.5,1.5);

  % Secondary condensate enters from the condenser and leaves as steam.
  \draw[->, very thick, codexblue]
    (9.5,0.85) -- (9.0,0.85)
    .. controls (8.5,0.85) and (8.5,1.5) .. (8.5,2.5);
  \draw[->, very thick, codexred]
    (8.5,2.5) .. controls (8.5,3.6) and (8.9,4.8) .. (9.5,5.0);

  \node[font=\tiny, align=center, text width=1.2cm]
    at (7.25,5.3) {Primary coolant};
  \node[font=\tiny, align=center, text width=1.4cm]
    at (8.75,5.3) {Secondary water/steam};
  % Gold arrows show heat transfer only; the two fluid loops do not mix.
  \foreach \y in {2.0,3.2,4.2} {
    \draw[->, thick, dashed, codexgold] (7.7,\y) -- (8.3,\y);
  }
  \node[codexgold, font=\tiny] at (8.0,4.45) {Heat};

  % Steam to turbine
  \draw[->, very thick, codexred!70]
    (9.5,5.0) -- (11.5,5.0);
  \node[above, font=\small] at (10.5,5.1) {Steam};

  % === Turbine ===
  \fill[gray!30] (11.5,4.0) rectangle (13.0,6.0);
  \draw[thick] (11.5,4.0) rectangle (13.0,6.0);
  \node[font=\tiny] at (12.25,5.55) {Turbine};

  % Generator
  \fill[gray!30] (13.0,4.0) rectangle (14.5,6.0);
  \draw[thick] (13.0,4.0) rectangle (14.5,6.0);
  \node[font=\tiny] at (13.75,4.45) {Generator};

  % Electricity out
  \draw[->, very thick, codexgold]
    (14.5,5.0) -- (15.5,5.0);
  \node[right, codexgold, font=\small] at (15.6,5.0)
    {Electricity};

  % Turbine exhaust steam enters the condenser.
  \draw[->, gray] (12.25,4.0) -- (12.25,3.0);
  \node[right, gray, font=\tiny] at (12.35,3.0)
    {Condenser};

  % Condensate returns from the condenser to the heat exchanger.
  \draw[->, codexblue!60]
    (12.25,3.0) -- (12.25,0.85) -- (9.5,0.85);
  \node[below, codexblue!60, font=\tiny]
    at (11.0,0.85) {Condensate return};

  % Shielding
  \draw[very thick, brown!70, dashed]
    (-0.5,-1.5) rectangle (5.0,7.5);
  \node[brown!70, font=\tiny] at (-0.3,7.8)
    {Concrete\\shielding};

\end{tikzpicture}
\end{center}
\end{diagbox}

% ============================================================
\section{Nuclear Fusion}
% ============================================================

\begin{definitionbox}[Nuclear Fusion]
\textbf{Nuclear fusion} is the joining of two light
nuclei to form a heavier nucleus, releasing energy.
Energy is released because the product has higher
binding energy per nucleon than the reactants.

The most promising fusion reactions for power
generation involve isotopes of hydrogen:
\[
  {}^2_1\text{H} + {}^3_1\text{H}
  \rightarrow {}^4_2\text{He} + {}^1_0\text{n}
  + 17.6\ \text{MeV}
\]
(Deuterium + Tritium $\rightarrow$ Helium-4
+ neutron + $17.6\ \text{MeV}$)

\textbf{Why fusion is difficult on Earth:}
Nuclei are positively charged and repel each other
strongly at normal temperatures. To overcome this
Coulomb barrier, temperatures of $\sim 10^8\ \text{K}$
are required (hotter than the Sun's core). At such
temperatures, matter exists as a \textbf{plasma} ---
all electrons stripped from nuclei. Confining a
plasma at $10^8\ \text{K}$ is the central engineering
challenge of fusion power.

\textbf{Advantages of fusion:}
\begin{itemize}[noitemsep]
  \item Fuel (deuterium) abundant in seawater
  \item No long-lived radioactive waste
  \item Much more energy per unit mass than fission
  \item No risk of runaway chain reaction
\end{itemize}
\end{definitionbox}

\begin{realworldbox}[Nuclear Power and Nigeria]
Nigeria does not currently operate any nuclear power
plants but has expressed long-term interest in
nuclear energy through the Nigerian Atomic Energy
Commission (NAEC). The argument for nuclear power
in Nigeria rests on its ability to deliver large,
reliable base-load power independent of weather
(unlike solar or wind) and without the carbon
emissions of natural gas. The argument against
centres on cost, safety infrastructure, long
construction times, and the availability of
abundant domestic gas reserves.

The physics of nuclear power in this chapter ---
binding energy, chain reactions, critical mass, and
waste decay --- is the technical foundation any
informed person needs to evaluate these arguments.
\end{realworldbox}

% ============================================================
\section{Radiation Safety and Biological Effects}
% ============================================================

\begin{defbox}[Radiation Dose and Safety]
\textbf{Absorbed dose} (gray, Gy):
energy deposited per unit mass of tissue.
$1\ \text{Gy} = 1\ \text{J\,kg}^{-1}$.

\textbf{Equivalent dose} (sievert, Sv):
absorbed dose weighted by the type of radiation
(alpha is $\sim 20\times$ more damaging per unit
energy than beta/gamma):
$H = wD$ where $w$ is the radiation weighting factor.

\textbf{Background radiation sources:}
Radon gas ($\sim 50\%$), medical X-rays ($\sim 15\%$),
cosmic rays, naturally radioactive materials in
food and ground. Annual UK/Nigerian average:
$\sim 2$--$3\ \text{mSv\,year}^{-1}$.

\textbf{Radiation protection principles:}
\begin{itemize}[noitemsep]
  \item \textbf{Distance:} Intensity falls as $1/r^2$.
  \item \textbf{Shielding:} Paper (alpha), aluminium
  (beta), lead/concrete (gamma/neutrons).
  \item \textbf{Time:} Minimise exposure time.
\end{itemize}
\end{defbox}

% ============================================================
\section{Chapter Summary}
% ============================================================

\begin{summarybox}
\textbf{Nucleus:} protons ($Z$) + neutrons ($N$);
$A = Z + N$. Nuclear radius $R = R_0 A^{1/3}$.

\textbf{Radioactive emissions:}
$\alpha$: ${}^4_2\text{He}$, $Z-2$, $A-4$;
$\beta^-$: electron, $Z+1$; $\beta^+$: positron, $Z-1$;
$\gamma$: photon, $Z$ and $A$ unchanged.

\textbf{Conservation:} $A$ and $Z$ conserved in all
nuclear equations.

\textbf{Decay is spontaneous and random.}

Activity: $A = \lambda N$;
$\lambda = \ln 2/t_{1/2}$.

\textbf{Exponential decay:}
$N = N_0 e^{-\lambda t}$; $A = A_0 e^{-\lambda t}$.

\textbf{Binding energy:}
$E_B = \Delta m \cdot c^2$; $1\ \text{u} = 931.5\ \text{MeV}$.
Peaks at ${}^{56}$Fe ($8.8\ \text{MeV/nucleon}$).

\textbf{Fission:} heavy nucleus splits;
$\sim 200\ \text{MeV}$ per event; chain reaction
(controlled in reactor, uncontrolled in bomb).

\textbf{Fusion:} light nuclei join; $10^8\ \text{K}$
required; huge energy per unit mass; clean fuel;
still under development for power.

\textbf{${}^{14}$C dating:} $t_{1/2} = 5730$ years;
for organic material up to $\sim 50{,}000$ years.
\end{summarybox}

% ============================================================
\newpage
\begin{exercisebox}[28]

\subsection*{Section A: Multiple Choice}

\textbf{A1.} An alpha particle has the same composition
as:

\mcoption{A}{A helium-3 nucleus}
\mcoption{B}{A helium-4 nucleus}
\mcoption{C}{Two protons only}
\mcoption{D}{A hydrogen-2 nucleus}
\ansref{Ch28-A1}

\vspace{0.4cm}

\textbf{A2.} In beta-minus ($\beta^-$) decay:

\mcoption{A}{A proton is emitted from the nucleus}
\mcoption{B}{The proton number $Z$ decreases by 1}
\mcoption{C}{A neutron converts to a proton and
an electron is emitted}
\mcoption{D}{The mass number $A$ decreases by 1}
\ansref{Ch28-A2}

\vspace{0.4cm}

\textbf{A3.} A radioactive sample has a half-life
of $20$ minutes. After $1$ hour, the fraction of
the original activity remaining is:

\mcoption{A}{$1/3$}
\mcoption{B}{$1/4$}
\mcoption{C}{$1/8$}
\mcoption{D}{$1/6$}
\ansref{Ch28-A3}

\vspace{0.4cm}

\textbf{A4.} Which radiation is most penetrating?

\mcoption{A}{Alpha}
\mcoption{B}{Beta-minus}
\mcoption{C}{Gamma}
\mcoption{D}{They all have equal penetration}
\ansref{Ch28-A4}

\vspace{0.4cm}

\textbf{A5.} The binding energy per nucleon is
greatest for:

\mcoption{A}{Hydrogen-1}
\mcoption{B}{Helium-4}
\mcoption{C}{Iron-56}
\mcoption{D}{Uranium-238}
\ansref{Ch28-A5}

\vspace{0.4cm}

\textbf{A6.} In a nuclear reactor, the function
of the moderator is to:

\mcoption{A}{Absorb excess neutrons to control
the chain reaction}
\mcoption{B}{Slow down fast neutrons to thermal
speeds}
\mcoption{C}{Cool the reactor core}
\mcoption{D}{Shield workers from radiation}
\ansref{Ch28-A6}

\vspace{0.4cm}

\textbf{A7.} Nuclear fusion releases energy because:

\mcoption{A}{The product nucleus is heavier than
the reactants}
\mcoption{B}{The product nucleus has lower binding
energy per nucleon than the reactants}
\mcoption{C}{The product has higher binding energy
per nucleon, so it has less mass-energy, and the
difference is released}
\mcoption{D}{Light nuclei repel each other,
releasing potential energy}
\ansref{Ch28-A7}

\vspace{0.4cm}

\textbf{A8.} A nucleus ${}^A_Z\text{X}$ undergoes
alpha decay followed by two beta-minus decays.
The final proton number is:

\mcoption{A}{$Z-2$}
\mcoption{B}{$Z-4$}
\mcoption{C}{$Z$}
\mcoption{D}{$Z+2$}
\ansref{Ch28-A8}

\vspace{0.4cm}

\textbf{A9.} The decay constant $\lambda$ of a
radioactive nuclide is related to its half-life
$t_{1/2}$ by:

\mcoption{A}{$\lambda = t_{1/2}/\ln 2$}
\mcoption{B}{$\lambda = \ln 2 \times t_{1/2}$}
\mcoption{C}{$\lambda = \ln 2 / t_{1/2}$}
\mcoption{D}{$\lambda = 1/t_{1/2}$}
\ansref{Ch28-A9}

\vspace{0.4cm}

\textbf{A10.} Radiocarbon dating uses carbon-14
because:

\mcoption{A}{It is the most abundant isotope of carbon}
\mcoption{B}{Its half-life of 5730 years is
appropriate for dating objects thousands of years old}
\mcoption{C}{It emits alpha particles, which are
easily detected}
\mcoption{D}{It is produced only in living organisms}
\ansref{Ch28-A10}

\vspace{0.6cm}

\subsection*{Section B: Theory and Calculations}

\textbf{B1.} Uranium-238 undergoes a decay series
to eventually reach stable lead-206.\\[4pt]
(a) Determine how many alpha decays and how many
beta-minus decays are involved. (Hint: write $A$
and $Z$ equations using the fact that each
$\alpha$ reduces $A$ by 4 and $Z$ by 2, while
each $\beta^-$ increases $Z$ by 1 with $A$
unchanged.)\\
(b) Write the nuclear equation for the first decay
in the series: ${}^{238}_{92}$U to ${}^{234}_{90}$Th.\\
(c) Thorium-234 then undergoes two consecutive
$\beta^-$ decays. Identify the nuclides produced
after each decay.\\
(d) The half-life of ${}^{238}$U is $4.47\times10^9$
years. Calculate the decay constant in s$^{-1}$.\\
(e) Calculate the activity of $1.00\ \text{g}$ of
${}^{238}$U (Avogadro's number $= 6.02\times10^{23}$
mol$^{-1}$, molar mass $= 238\ \text{g\,mol}^{-1}$).
\hfill\ansref{Ch28-B1}

\vspace{0.4cm}

\textbf{B2.} A sample of radioactive iodine-131
(half-life $= 8.04$ days) has an initial activity
of $3.7\times10^9\ \text{Bq}$.
($N_A = 6.02\times10^{23}\ \text{mol}^{-1}$,
molar mass of I-131 = 131 g\,mol$^{-1}$)\\[4pt]
(a) Calculate the decay constant in s$^{-1}$.\\
(b) Calculate the initial number of ${}^{131}$I
atoms.\\
(c) Calculate the initial mass of ${}^{131}$I.\\
(d) Calculate the activity after $32.16$ days.\\
(e) Plot (sketch) the decay curve of activity vs
time, labelling the axes and marking the values
at $t = 0$, $t_{1/2}$, $2t_{1/2}$, $3t_{1/2}$.\\
(f) Iodine-131 is used for thyroid treatment. A
patient needs to receive an effective dose over
$32$ days. Why does the doctor arrange for the
patient to receive the dose immediately rather
than waiting to administer it later?
\hfill\ansref{Ch28-B2}

\vspace{0.4cm}

\textbf{B3.} (a) Calculate the binding energy of
${}^{56}_{26}$Fe (iron-56), often called the most
stable nucleus.\\
Given: $m({}^{56}$Fe$) = 55.934939\ \text{u}$;
$m_p = 1.007276\ \text{u}$;
$m_n = 1.008665\ \text{u}$;
$1\ \text{u} = 931.5\ \text{MeV}/c^2$.\\[4pt]
(b) Calculate the binding energy per nucleon.\\
(c) Compare with ${}^4_2$He
($E_B = 28.3\ \text{MeV}$, $A = 4$).
Which has greater binding energy per nucleon?\\
(d) Use the binding energy curve (Figure in chapter)
to explain qualitatively:\\
(i) Why nuclear fission of ${}^{235}$U releases energy.\\
(ii) Why nuclear fusion of hydrogen releases energy.\\
(iii) Why neither fission of ${}^{56}$Fe nor fusion
of ${}^{56}$Fe releases energy.
\hfill\ansref{Ch28-B3}

\vspace{0.4cm}

\textbf{B4.} An archaeologist finds a piece of
ancient wood. She measures the activity of the
carbon-14 in $10\ \text{g}$ of the wood sample
to be $1.5\ \text{Bq}$.
A living tree gives $6.5\ \text{Bq}$ per $10\ \text{g}$
of carbon.
($t_{1/2}$ of ${}^{14}$C $= 5730$ years,
$\ln 2 = 0.693$)\\[4pt]
(a) Calculate the fraction of ${}^{14}$C remaining.\\
(b) Calculate the age of the wood.\\
(c) The measurement uncertainty in the activity
is $\pm 0.2\ \text{Bq}$. Estimate the uncertainty
in the calculated age.\\
(d) Explain why radiocarbon dating is not reliable
for samples older than about 50,000 years.\\
(e) The sample is from an ancient Yoruba wooden
carving. If the age calculated is $1200$ years,
place this in historical context (what period of
Yoruba history might this correspond to?).
\hfill\ansref{Ch28-B4}

\vspace{0.4cm}

\textbf{B5.} A nuclear power reactor produces
$1000\ \text{MW}$ of electrical power at an
efficiency of $35\%$.
(Energy per fission of ${}^{235}$U $= 200\ \text{MeV}$;
$1\ \text{MeV} = 1.6\times10^{-13}\ \text{J}$;
$N_A = 6.02\times10^{23}\ \text{mol}^{-1}$;
molar mass of ${}^{235}$U $= 235\ \text{g\,mol}^{-1}$)\\[4pt]
(a) Calculate the thermal power (heat power) that
the reactor must produce.\\
(b) Calculate the number of fissions per second.\\
(c) Calculate the mass of ${}^{235}$U consumed per
day.\\
(d) Compare the mass of ${}^{235}$U needed to power
this reactor for one year with the mass of coal
needed to produce the same electrical energy from
a coal power station (efficiency $40\%$, energy
density of coal $\approx 24\ \text{MJ\,kg}^{-1}$).\\
(e) One concern about nuclear power is the
disposal of radioactive waste. A spent fuel rod
contains fission fragments including ${}^{90}$Sr
(half-life $= 29$ years). Calculate how many
half-lives must pass before the ${}^{90}$Sr
activity falls to $0.1\%$ of its initial value.
How many years is this?
\hfill\ansref{Ch28-B5}

\end{exercisebox}
